\documentclass{resume} % Use the custom resume.cls style

\usepackage[left=0.4 in,top=0.4in,right=0.4 in,bottom=0.4in]{geometry} % Document margins
\usepackage{enumitem}
\usepackage{marvosym}

\newcommand{\tab}[1]{\hspace{.2667\textwidth}\rlap{#1}} 
\newcommand{\itab}[1]{\hspace{0em}\rlap{#1}}
\setlength{\itemsep}{0pt}

\name{Erle Zhu} % Your name
% You can merge both of these into a single line, if you do not have a website.
\address{Beijing, China} 
\address{\href{mailto:erlerzhu@gmail.com}{Gmail} \\  
\href{https://github.com/Lez-3f}{Github} \\
\href{ https://scholar.google.com/citations?user=z5GT6QsAAAAJ&hl=zh-CN} {Scholar}
}  %
% \href{https://linkedin.com/company/faangpath}{linkedin.com/company/faangpath} \\
\begin{document}

%----------------------------------------------------------------------------------------
%	OBJECTIVE
%----------------------------------------------------------------------------------------

% \begin{rSection}{OBJECTIVE}

% {Software Engineer with 2+ years of experience in XXX, seeking full-time XXX roles.}

%----------------------------------------------------------------------------------------
%	EDUCATION SECTION
%----------------------------------------------------------------------------------------

\begin{rSection}{Education}

{\bf PhD in Computer Science} \hfill {2024-2029 (Expected)} \\
Department of Computer Science \& Technology, Tsinghua University \hfill \textit{Beijing, China}
\begin{itemize}
    \item Member of CoAI research group
    \item Research Goal: Building generalizable foundation agent via Reinforcement Learning. 
\end{itemize}

{\bf Bachelor of Mathematics and Physics + EE} \hfill {2020-2024}\\
Weiyang College, Tsinghua University \hfill \textit{Beijing, China}
\begin{itemize}
    \item Major GPA: 3.93/4.00 (Rank: 2/30)
    \item Got full grade(4.0/4.0) in All courses related to Mathematics, Physics and Data Science.
\end{itemize}
%Minor in Linguistics \smallskip \\
%Member of Eta Kappa Nu \\
%Member of Upsilon Pi Epsilon \\


\end{rSection}

%----------------------------------------------------------------------------------------
% EXPERIENCE
%----------------------------------------------------------------------------------------
\begin{rSection}{EXPERIENCE}

\textbf{Z.AI} \\
Research Intern \hfill February 2025 - Precent
 \begin{itemize}
    \itemsep -3pt {} 
     \item Data Synthesis \& RLHF for GLM4.5's Instruction Following. 
     \item Optimized the RL infra \textit{Slime} for Agentic tasks, supporting multi-task Agentic RL. 
     \item Optimized GLM5 series in Terminal-Agent Scenarios, involves data synthesis, SFT, RLVR. 
\end{itemize}

\textbf{LeapLab, Tsinghua University} \hfill \textit{Beijing, China} \\
Research Assistant \hfill July 2022 - March 2023
 \begin{itemize}
    \itemsep -3pt {} 
     % \item Advisor: \href{http://www.gaohuang.net/}{Gao Huang}
     \item Researched on how to augmenting the exploration of RL Agent with unsupervised reward signal, introduced a retrieval augmented technique into unsupervised RL setting and gained SoTA in URL-Benchmark. 
     % \item Paper \textit{"Augmenting Unsupervised Reinforcement Learning with Self-Reference"} in submission to ICML2023 \\(co-first author)
 \end{itemize}

\textbf{Energy Internet Research Institute, Tsinghua University} \hfill \textit{Sichuan, China}\\
Research Intern \hfill June 2023 - July 2023
\begin{itemize}
    \itemsep -3pt {} 
     % \item Advisor: \href{http://www.gaohuang.net/}{Gao Huang}
     \item Worked on data engineering \& machine learning for energy consumption prediction.
\end{itemize}
\end{rSection}

%----------------------------------------------------------------------------------------
%	pub
%----------------------------------------------------------------------------------------

\begin{rSection}{publications}

% Vew all my publications in my \href{https://scholar.google.com/citations?user=z5GT6QsAAAAJ&hl=zh-CN}{scholar page}.

\begin{enumerate}[label={[\arabic*]}]
    \item \textbf{Data Efficient RLVR via Off-Policy Influence Guidance} \\
        \textbf{Erle Zhu*}, Dazhi Jiang*, Yuan Wang, Xujun Li, Jiale Cheng, Yuxian Gu, Yilin Niu, Aohan Zeng, Jie Tang, Minlie Huang, Hongning Wang\textsuperscript{\Letter} \\
        \textit{ACL 2026 (Main)} 
    \item \textbf{MAPS: Advancing Multi-Modal Reasoning in Expert-level Physical Science} \\
        \textbf{Erle Zhu}, Yadi Liu, Zhe Zhang, Xujun Li, JinZhou, Xinjie Yu, Minlie Huang, Hongning Wang\textsuperscript{\Letter} \\
        \textit{ICLR, 2025.}
    \item \textbf{Augmenting Unsupervised Reinforcement Learning with Self-Reference} \\
        Andrew Zhao*, \textbf{Erle Zhu*}, Rui Lu, Matthieu Gaetan Lin, Ning Jia, Yong-jin Liu, Gao Huang\textsuperscript{\Letter} \\
        \textit{Neural Network, 2025.}

    % \item \textbf{AI teaching assistants help the collaborative exploration of large-scale individualized teaching of Circuit Principle course} \\
    %    Pengcheng Zhang, Yi Ren, Xinjie Yu\textsuperscript{\Letter},  \textbf{Erle Zhu}\\
    %    The 14th Annual Conference of Circuits and Signal Systems, Electromagnetic Field Teaching and Materials Research Association of Chinese Colleges, 2024(Language: Chinese, \textbf{Outstanding Paper Reward})
    %    % 之后仅保留在中文简历中：
    %    % Al助教助力“电路原理”规模化因材施教的协作探索
    %    % 章鹏程、于歆杰、任意、祝尔乐
    %    % 高等学校电路和信号系统、电磁场教学与教材研究会第十四届年会
    %    % 优秀论文奖
       
\end{enumerate}
    (* indicates equal contribution)
\end{rSection} 

%----------------------------------------------------------------------------------------
%	projects
%----------------------------------------------------------------------------------------

\begin{rSection}{PROJECTS}
View all my projects in \href{https://github.com/Lez-3f/}{my Github}. \\
\vspace{-1.25em}
\item \href{https://github.com/zai-org/GLM-5} {\textbf{GLM5 Series} }
\\{A top-tier foundation model for coding and agentic tasks. 
I contributed to the Agentic RL \& Data Synthesis in Terminal Agent}
\item \href{https://github.com/zai-org/GLM-4.5} {\textbf{GLM4.5 Series} }
\\{A top-tier foundation model for coding and agentic tasks. 
I contributed to the Data Pipeline of RLHF and RL Infra for multi-task Agentic RL. }
\item \href{https://github.com/THUDM/slime} {\textbf{Slime} }
\\{An open-source post-traning framework and the training infra for GLM. 
I contributed to the fully asynchronous RL training implementation in open-source repo \& internal version in Zhipu for agentic RL.}
\item \href{https://yuketang.tsinghua.edu.cn/ai} { \textbf{THU-AI4Edu} }
\\{An AI-driven education reform project in Tsinghua University. I led the Educational Agent project for circuit theory \& programming in Tsinghua University.}
\item \href{https://github.com/GJCav/thuwy} { \textbf{WeiWeiYang} }
\\{A Wechat program aiming to facilitate the daily lives of students in Weiyang College. 
% It includes the functions of public facilities rental, message notification, venue reservation, academic Q\&A and so on. 
I contributed to the system backend for email and message notification support. }
\item \href{https://github.com/Lez-3f/THU-EDC2021-BugCar} {\textbf{EDC-BugCar} }
\\{A Project of the 23th \href{https://mp.weixin.qq.com/s/ywimZvBIiEm3Ct4eLW689Q}{THU-EDC}(Top hardware competition in THU). Our team won the third prize (top 8 teams) out of 56 teams. I contributed to the control algorithm.  }

\end{rSection} 

%----------------------------------------------------------------------------------------
% TECHINICAL STRENGTHS	
%----------------------------------------------------------------------------------------
\begin{rSection}{SKILLS}

\begin{tabular}{ @{} >{\bfseries}l @{\hspace{6ex}} l }
LLM Framework: Slime, VeRL, OpenRLHF, SGLang\\
Software: Python (Pytorch, Tensorflow), C/C++, MATLAB (Simulink), \LaTeX, MySQL \\
Hardware: MCU(STM32, MSP430, Arduino), FPGA(Verilog)
\end{tabular}\\
\end{rSection}

%----------------------------------------------------------------------------------------
% Honors
%----------------------------------------------------------------------------------------
\begin{rSection}{HONORS}

\begin{itemize}
    \item Comprehensive-Merit Scholarship of THU, 2021, 2022 (Top 15\%)
    \item Technology-Innovation Scholarship of THU, 2022 (Top 10\%)
    \item Academic-Merit Scholarship of THU, 2021 (First order, Top 3\%)
    \item The 3rd Prize of 23th THU-EDC(Electronic Design Competition), 2021 (8/56)
    \item The 1st Prize of the 13th Chinese Mathematics Competition for College Students, 2021
\end{itemize}
\end{rSection}

%----------------------------------------------------------------------------------------
% Service
%----------------------------------------------------------------------------------------
\begin{rSection}{Services}

% \begin{itemize}
%     \item Minister @WY-SAST-Academic(Academic dept., Student Association of Science and Technology) (2022.2-2023.2)
% \end{itemize}

\textbf{Minister} \hfill Feb 2022 - Feb 2023\\
Academic Dept., Student Association of Science \& Technology \hfill \textit{Weiyang College, Tsinghua Univ.}
\begin{itemize}
    \item Holding lectures of academic skills for undergraduates
\end{itemize}

\textbf{Teaching Assistant} \\
I have been a TA for following courses:
\begin{itemize}
    \item LLM Summer School 2024
    \item Linear Algebra (English) 2024 Fall
\end{itemize}
\textbf{Reviewer} \\
I have been a reviewer for 
\begin{itemize}
    \item NeurIPS(2024, 2025)
\end{itemize}

\end{rSection}

\textbf{Talks \& Lectures} \\
I gave the following lectures or talks:
\begin{itemize}
    \item Tutorials of Weiyang City Competition  - Lecture 2 (Machine Learning)\hfill \textit{Weiyang College, Tsinghua Univ.}
\end{itemize}


% \begin{rSection}{Misc}

% % \begin{itemize}
% %     \item Minister @WY-SAST-Academic(Academic dept., Student Association of Science and Technology) (2022.2-2023.2)
% % \end{itemize}

% \textbf{Minister} \hfill Feb 2022 - Feb 2023\\
% Academic Dept., Student Association of Science \& Technology \hfill \textit{Weiyang College, Tsinghua Univ.}
% \begin{itemize}
%     \item Holding lectures of academic skills for undergraduates
% \end{itemize}

% \end{rSection}


\end{document}
